% ===========================================================================
\clearpage
\floraon
\florathispage{layout=surround, seed=41, density=0.32, visitors=0.08}
\part{Evidence and operation}
\thispagestyle{empty}
\clearpage
\floraoff
% ===========================================================================

\section{Testing, CI, and reproducibility}
\label{sec:repro}

\subsection{The four test layers}
\label{sec:test-layers}

The gate is layered so a regression stops at the cheapest test that can see
it:
\begin{enumerate}
  \item \textbf{gtest / ctest.} C++ unit and integration tests carry one
        hardware label: \kst{no-gpu}, \kst{gpu-amd}, or
        \kst{gpu-nvidia-rental}. The first two are the local gate; the third is
        owner-executed on a rental.
  \item \textbf{lit / FileCheck.} The \kst{check-kernscope} target runs the
        parser, pass, verifier-manifest, and CLI tests against committed IR and
        ISA corpora.
  \item \textbf{NumPy bf16 golden harness.} \kst{pytest tests/golden} compares
        generated kernels against reference outputs under the bf16 tolerance
        contract. It is the independent implementation of the C++ golden gate.
  \item \textbf{Generated-code compile checks.} \kst{scripts/compile_checks.sh}
        compiles the template variants with \kst{nvcc} for sm\_80/sm\_90
        (compile-only) and \kst{hipcc} for gfx1101, then runs the live HIP
        variant through the golden gate.
\end{enumerate}

\subsection{The CI gate}
\label{sec:ci}

\kst{scripts/ci.sh} runs the stages in order and exits non-zero on the first
failure. A stage whose artifact is absent prints \kst{[SKIP]} and stays green,
which is a bootstrap convenience, not a claim that the stage ran.

\begin{table}[htbp]
\centering
\small
\begin{tabularx}{\textwidth}{@{}l X@{}}
\toprule
Stage & Check \\
\midrule
a & \kst{scripts/configure.sh} and \kst{cmake --build build}. \\
b & \kst{ctest -L "no-gpu|gpu-amd"} with hardware labels. \\
b2 & Non-decreasing count guard against \kst{scripts/ci-baseline.txt} (550). \\
c & \kst{cmake --build build --target check-kernscope} (lit/FileCheck). \\
d & \kst{pytest tests/golden -q} (NumPy bf16 goldens). \\
e & \kst{scripts/compile_checks.sh} (nvcc sm\_80/90, hipcc gfx1101, live HIP golden). \\
f & Committed fixtures total at most 64\,MB (\kst{du -sm tests/fixtures}). \\
g & NVIDIA-header guards: no vendored CUDA/CUPTI/NVTX header filenames or vendor copyright markers outside the documented exceptions. \\
\bottomrule
\end{tabularx}
\caption{The CI stages. The count guard is non-decreasing, not exact: adding
tests is always green, and removing a registration below the floor fails.}
\label{tab:ci}
\end{table}

The \kst{CI_SKIP_CONFIGURE=1} escape hatch skips configure and build. It exists
for a bootstrap situation where the nix-store LLVM resolution is being
repaired; it is owner-set and it makes the skip visible. The gate does not run
the dashboard validator, and it does not run a secret scanner. Those are
operational gaps, not hidden passes.

\subsection{The master reproducibility script}
\label{sec:reproduce-all}

\kst{scripts/reproduce-all.sh} chains four stages: the lock check,
\kst{ci.sh}, \kst{demo.sh}, and the dashboard. The lock check requires
\kst{devenv.lock} to exist and be tracked; it never regenerates the lock
silently. The dashboard writes \fpath{reports/dashboard.html} and
\fpath{reports/dashboard.json} with \kst{SOURCE_DATE_EPOCH} defaulted to the
HEAD commit time, so two runs on the same commit can be byte-identical.

A clean checkout is:
\begin{lstlisting}
devenv shell
git clone https://github.com/meillaya/kernscope.git
cd kernscope
bash scripts/ci.sh
bash scripts/reproduce-all.sh
\end{lstlisting}
The first configure and build are the slowest (roughly ten minutes on the
reference machine); later runs reuse \fpath{build/}. The committed fixtures,
golden inputs, lit corpora, and schema files are the complete gate input. The
earlier documentation used a local absolute clone path; this report states the
public URL instead.

\subsection{Owner-gated NVIDIA runs}
\label{sec:rental-runbook}

Only \kst{scripts/rental-validate.sh} runs NVIDIA code. It is never called by
\kst{ci.sh}. In \kst{--dry-run} it prints the sequence: SSH to the rental host,
check \kst{nvcc}, detect the CUPTI version from \kst{cupti_version.h}, run
\kst{ctest -L gpu-nvidia-rental}, and place evidence under
\fpath{reports/rentals/}. In exec mode it refuses when \kst{CI=true}, requires
a configured build tree, probes \kst{$CUDA_HOME}, \fpath{/usr/local/cuda}, and
\fpath{/opt/cuda} for CUPTI, and writes a timestamped summary with the verdict,
exit code, CUDA major, and CUPTI version.

The gate is procedural. The script checks the environment, not the operator's
identity, and it assumes SSH keys are already configured. The repository
ignores \fpath{reports/rentals/*} except the placeholder, so an owner who wants
to commit a rental summary must force-add it. No rental numbers are claimed in
this report.

\subsection{Secrets and sensitive data}
\label{sec:secrets}

The repository ignores \fpath{private_key.pem}, \fpath{public_key.pem},
\fpath{.env}, and \fpath{.env.*} (with an exception for \fpath{.env.example}).
The rental script does not manage keys. There is no secret scan in \kst{ci.sh};
the rule is enforced by the ignore list and by review, not by an automated
gate. Capture bundles are local-only and gitignored; snapshot sizes are capped
at 1\,GiB by default, and the fixture budget is 64\,MB.

% ===========================================================================
\section{Evaluation}
\label{sec:eval}

\subsection{The demo run}
\label{sec:eval-demo}

\kst{scripts/demo.sh} ran the loop on the local gfx1101 in 153\,s against a
900\,s budget. The aggregated report is the primary cross-tool evidence.

\begin{table}[htbp]
\centering
\small
\begin{tabularx}{\textwidth}{@{}l l X@{}}
\toprule
Stage & Status & Evidence \\
\midrule
1 profile & ok & 3 kernels, 19.6\,$\mu$s wall, hot layernorm 9.2\,$\mu$s, 24 registers, occupancy 1, \kst{memory_bound} hint (AI 1.375, ridge 27.614). \\
2 capture & ok & Bundle \fpath{bundle.v0}, 3 blobs, grid $1024\times1\times1$, block $128\times1\times1$, bitwise replay PASS, reference SHA-256 \kst{a0027e...49eec}. \\
3 analyze & ok & Captured AMDGPU ISA, 216 ALU ops, \kst{possibly_divergent_branch} warning, three informational findings. \\
4 capture-tune & ok & 4 candidates validated, baseline block 128 at 0.01208\,ms, best block 256 at 0.01206\,ms, speedup $1.0017\times$. \\
5 verify + pass & ok & 13 bias\_gelu variants, 12 compiled, 1 gated; best 0.17184\,ms vs 0.20332\,ms, $1.1832\times$; pass vector width 1$\rightarrow$4, loads 0$\rightarrow$12, uncoalesced 64$\rightarrow$16. \\
\bottomrule
\end{tabularx}
\caption{The five demo stages. The report's cross-tool block links the stage-3
analyzer finding ids to the stage-1 profiler hint.}
\label{tab:demo}
\end{table}

The cross-tool link is the part that matters. The profiler's
\kst{memory_bound} hint carries \kst{analyzer_cross_refs} with the finding ids
\kst{possibly_divergent_branch}, \kst{no_tensor_core_usage},
\kst{occupancy_not_applicable}, and \kst{suggestions_not_applicable}, and the
stage-3 evidence path. The analyzer ran on the captured kernel's
disassembled ISA, not on a hand-picked example. The verifier removed one
variant before compilation. The pass delta used the analyzer's own metrics.
This is the concrete meaning of ``the loop closes'' in this report.

\subsection{Analyzer and verifier evidence}
\label{sec:eval-static}

The analyzer golden fixture is described in \Cref{sec:c1-evidence}. The
verifier corpus has 34 entries across 19 source files, with 18 \kst{SAFE}
verdicts and 16 non-\kst{SAFE} verdicts spread across the bounds, race,
barrier, pointer, uninitialized-read, and address-space classes. Both results
are \kst{projected}: they come from static analysis, not from a GPU run. The
verifier's anti-self-report gate diffs the CLI's JSON against the manifest read
from disk, so a corrupted expectation fails the lit suite.

\subsection{Tuner and capture evidence}
\label{sec:eval-live}

The capture-tune and verify-tune results in \Cref{tab:demo} are measured on
gfx1101. The capture path is bitwise: every explored config must reproduce the
captured output hash before it is timed. The template path is
correctness-gated: a variant that fails cosine, relative error, or PCC is
discarded, not timed. The measured speedups are small for capture-tune
(1.0017$\times$) and modest for the bias\_gelu template search
(1.1832$\times$); the report quotes them as measured, not as a general
speedup claim.

The dashboard records a separate matmul search with 15 measured variants and a
best median of 415.324\,$\mu$s. The raw JSONL for that run is not committed,
so this is a dashboard-sourced value, not a value reconstructed from a raw
record. The distinction matters because the dashboard's tune section lists
source paths that are absent from the tree.

\subsection{Pass and profile evidence}
\label{sec:eval-projected}

The layernorm pass delta is projected: vector width 1$\rightarrow$4,
vectorized global loads 0$\rightarrow$12, uncoalesced accesses
64$\rightarrow$16, registers unchanged at 175. The before/after PTX compiles
cleanly, and gemm and softmax are correctly no-transform. The report does not
claim a wall-clock pass speedup; it claims the static delta that the pipeline
measures.

The profiler evidence is measured: 3 kernels, 19.6\,$\mu$s wall, a
9.2\,$\mu$s hot kernel, 24 registers, occupancy 1, and a memory-bound hint
with arithmetic intensity 1.375 and ridge point 27.614. The hint is a
compiler-aware diagnosis from activity records and the roofline, not from
hardware counters.

\subsection{Threats to validity}
\label{sec:threats}

\begin{itemize}
  \item \textbf{One measured platform.} Every live number is gfx1101 on one
        RX 7800 XT. The same kernel on CDNA or NVIDIA may behave differently.
  \item \textbf{Unlocked clocks.} The peak run recorded clocks with
        \kst{rocm-smi} but did not lock them, and the machine was shared with
        development activity. The peaks are a description of this machine at
        that time.
  \item \textbf{Projected deltas.} The pass numbers are static IR/PTX
        artifacts. They are not wall-clock speedups.
  \item \textbf{Verifier scope.} The verifier is a may-analysis, not a formal
        prover. Its false-positive and false-negative classes are stated in
        \Cref{sec:c3-soundness}.
  \item \textbf{Dashboard gaps.} The tune section's source files are absent;
        the matmul number is quoted from the dashboard, and the validator is
        not part of \kst{ci.sh}. The report says so rather than presenting the
        dashboard as a complete raw archive.
  \item \textbf{Documentation consolidation.} The raw Markdown notes and the
        committed JSON/HTML evidence were folded into this report and removed
        from the tree. The values remain; the raw files do not. Regeneration is
        possible with the committed scripts on a gfx1101 machine.
\end{itemize}

% ===========================================================================
\section{Context and prior art}
\label{sec:prior}

kernscope is an engineering synthesis, not a new verification calculus. Its
design choices follow a body of work that the project's research briefs
survey; this section states the main lineages and what kernscope borrows or
declines.

\subsection{Performance models and static analysis}
\label{sec:prior-analysis}

The roofline model \citep{roofline2009} supplies the compute and memory
ceilings and the ridge point used by the tuner and profiler. Occupancy follows
the CUDA programming guide's resource model \citep{cudaProgrammingGuide}, with
the low-occupancy/high-ILP trade-off documented by Volkov
\citep{volkov2010}. Nsight Compute's Speed of Light vocabulary
\citep{nsightComputeSOL} is the diagnostic register the profiler adopts.

The analyzer's static estimates sit against a line of microbenchmark-driven
architecture studies (Volta, Turing, Hopper, and CDNA3) and against simulators
such as GPGPU-Sim and Accel-Sim \citep{gpgpusim2009,accelsim2020} and binary
instrumentation such as NVBit \citep{nvbit2019}. Those tools measure or
simulate; kernscope reads the ISA text and the architecture database, which is
cheaper and less precise. The analyzer's contribution is the reuse of the same
architecture constants and metrics by the tuner, pass report, and profiler.

\subsection{Autotuning and tensor compilation}
\label{sec:prior-tune}

The autotuner follows the algorithm/schedule separation of Halide
\citep{halide2013}, the search-based scheduling of TVM and Ansor
\citep{tvm2018,ansor2020}, and the tile-level autotuning model of Triton
\citep{triton2019}. CUTLASS \citep{cutlass} and the vendor libraries set the
performance bar; Kernel Tuner \citep{kerneltuner2024} is the closest standalone
autotuner prior art. FlashAttention \citep{flashattention2022} motivates the
memory-bound attention family and the IO-aware tiling direction. kernscope
does not generate a general tensor program; it generates family-specific
templates, gates them on correctness, and records the search in a cache that
can be replayed.

\subsection{GPU verification}
\label{sec:prior-verify}

GPUVerify \citep{gpuverify2012,gpuverify2015} is the methodological anchor:
the two-thread reduction turns race and divergence checks into sequential
program verification. PUG, SESA, GKLEE, and ESBMC-GPU
\citep{pug2010,sesa2014,gklee2012,esbmcgpu2017} explore symbolic execution and
bounded model checking. Faial \citep{faial2024} contributes behavioral types
over index expressions and the finding that most data-race-free kernels need a
thread-configuration assumption. HiRace and HGRD
\citep{hirace2024,hgrd2026} bring dynamic and MLIR-level analyses. LLOV
\citep{llov2020} demonstrates static race freedom on LLVM IR. The Linux eBPF
verifier \citep{ebpfVerifier} supplies the interaction model: typed ranges,
bounded state exploration, and a concrete rejection reason. MemorySanitizer
\citep{msan2015} and the SPIR-V validator \citep{spirvSpec} inform the
uninitialized-read and well-formedness checks.

kernscope-verify is deliberately smaller: eight may-analyses over LLVM IR at
\kst{-O1}, with the launch configuration ingested from the host half and a
three-tier exact/upper-bound/symbolic contract. It does not claim the
soundness of GPUVerify or Faial, and its false-positive and false-negative
classes are published rather than hidden.

\subsection{Capture and replay}
\label{sec:prior-capture}

Kerncap \citep{kerncap2026} is the closest complete system: HSA-level
extraction and VA-faithful replay for HIP and Triton. Mneme \citep{mneme2023}
records and replays kernels for tuning; the LLNL Record-Remix-Replay line
\citep{r3-2026} and Tally \citep{tally2024} explore the same design space.
CUPTI \citep{cupti} and rocprofiler-sdk \citep{rocprofilerSDK} provide the
vendor tracing surfaces. kernscope's contribution is the shared bundle model
across CUPTI, LD\_PRELOAD, and rocprofiler-sdk, the explicit patch table for
embedded pointers, and the use of the same bundle by the analyzer,
capture-tune, and the standalone reproducer.

\subsection{LLVM passes and profiling}
\label{sec:prior-pass}

The pass plugin uses LLVM's New Pass Manager and the out-of-tree
\kst{PassPluginLibraryInfo} ABI \citep{llvm2004,mlirCGO2021}. The
transformations are deliberately narrow: promotion, tiling, vectorization,
padding, redundant-load elimination, and a coalescing analysis. The GPU
backends already run address-space inference, load/store vectorization, and
redundant-load elimination; the project's value is the analyzer-measured
before/after delta, not a claim that these transformations are new. IREE
\citep{iree} is the closest production analog for a compiler-integrated
memory-optimization pipeline.

The profiler follows the activity-only tier of CUPTI and rocprofiler-sdk plus
the roofline and occupancy vocabulary of Nsight Compute
\citep{nsightComputeSOL}. It reports what it can evidence and degrades to
\kst{insufficient_evidence} when the family, PTX, or occupancy is unknown.

% ===========================================================================
\section{Limitations and future work}
\label{sec:limits}

\subsection{Limitations}
\label{sec:limitations-list}

\begin{itemize}
  \item \textbf{No NVIDIA runtime evidence.} CUPTI capture, CUPTI profiling,
        and NVIDIA runtime benchmarks are compile- and unit-tested only. The
        rental path exists but produced no committed numbers.
  \item \textbf{One measured GPU.} The measured results are gfx1101 only. The
        analyzer supports other architectures, but their occupancy and
        roofline numbers are not validated here.
  \item \textbf{Pass deltas are projected.} The vector-width and coalescing
        improvements are static IR/PTX measurements. A wall-clock speedup for
        the transformed layernorm is not claimed.
  \item \textbf{The verifier is a may-analysis.} Data-dependent indices and
        runtime allocations degrade to non-\kst{SAFE} verdicts; atomics,
        warp-synchronous code, and inter-block races are out of scope.
  \item \textbf{Capture format drift.} The v1 schema is the reading contract,
        but the demo bundle is v0 and the current record path still uses the
        v0 environment-variable flow. The v1 writer path is not exercised by
        the demo.
  \item \textbf{Dashboard incompleteness.} The committed dashboard's tune
        section cites per-run JSONL files that are absent, and the dashboard
        validator is not invoked by \kst{ci.sh}. The matmul number is quoted
        from the dashboard only.
  \item \textbf{Report schema does not enforce provenance.} The schema
        requires the kernel, metrics, limiter, warnings, and suggestions; the
        provenance object is a report-engine and dashboard convention.
  \item \textbf{Tooling gaps.} The auxiliary PTX dump source has no CMake
        target and is not built. The capture CLI's usage exit codes do not
        consistently match the documentation. The report names these rather
        than presenting a uniform interface that does not exist.
  \item \textbf{The archive is textual.} The consolidation removed the raw
        JSON/HTML evidence. The numbers in this report are preserved, but a
        reader who needs the raw records must regenerate them on a gfx1101
        machine with the committed scripts.
\end{itemize}

\subsection{Future work}
\label{sec:future}

\begin{enumerate}
  \item \textbf{Run the NVIDIA rental.} Execute \kst{rental-validate.sh} on an
        A100/H100 host, commit the summary with the same provenance schema, and
        add an NVIDIA column to the dashboard.
  \item \textbf{Exercise CUDA 13.} Build the record adapter against
        \kst{CUPTI_API_VERSION} 130001 and turn the Kernel9/Kernel10
        \kst{static_assert}s into a tested migration.
  \item \textbf{Close the v0/v1 capture gap.} Make \kst{record} write the v1
        bundle by default, exercise the schema in the demo, and remove the
        environment-variable path.
  \item \textbf{Strengthen the verifier.} Add HGRD-style barrier intervals,
        Faial-style index types, and a PTX well-formedness layer; keep the
        false-positive and false-negative tables current.
  \item \textbf{Wire evidence into CI.} Run the dashboard validator in
        \kst{ci.sh}, fix its failure path, and add a secret scan.
  \item \textbf{Roofline-driven pass selection.} Choose vectorization or tiling
        from the profiler's bound class and the pass report's predicted delta,
        then validate the choice with a measured benchmark.
  \item \textbf{Validate generated kernels.} Use the verifier and golden gate
        as a correctness harness for LLM-generated CUDA, with a kernel-benchmark
        sweep rather than a single family.
\end{enumerate}

% ===========================================================================
\section{Conclusion}
\label{sec:conclusion}

kernscope's claim is modest and testable: a GPU-kernel toolchain can share one
library stack, one report schema, and one provenance discipline, and its tools
can validate each other instead of self-reporting. The loop runs on one local
GPU in 153\,s. The analyzer reads the captured ISA, the verifier removes a
risky variant before compilation, the capture replay is bitwise, and the pass
delta is measured with the same metrics the analyzer uses. The numbers are
small where the work is small (a 1.0017$\times$ capture-tune speedup) and
larger where the search has room (a 1.1832$\times$ bias\_gelu speedup); the
report keeps them in their provenance categories.

The same discipline that makes the evidence useful makes the limits visible.
There are no NVIDIA runtime numbers, the pass deltas are projected, the
verifier is not a prover, one GPU supplied every live measurement, and the
dashboard is not a complete raw archive. Those are not caveats appended to the
claim; they are the conditions under which the claim is true. The repository's
scripts can regenerate the evidence on a gfx1101 machine, and this report is
the archival account of what was present when the documentation was
consolidated.

% ===========================================================================
